\section{Basic observables and the kernel $e$-projections}
\label{eproj:section}
% BEGIN SUBJECT INDEX: eproj:section
\index{kernel!idempotents}
\index{kernel!partition marginals}
\index{kernel!range marginals}
\index{transversal}
% END SUBJECT INDEX

Let $K$ be the minimum-rank ideal of a finite transformation semigroup,
with rank $r$, and let $\lambda$ be its limiting kernel measure. Denote
its partition and range marginals by $\alpha_i$ and $\beta_j$. Write
$e_{ij}$ for the idempotent with partition $\mathcal P_i$ and range
$\mathcal R_j$. Each range is a transversal of every kernel partition.
All matrices below use right-action conventions.

\subsection{The two basic expectation operators}
\label{notation:expectations}
% BEGIN SUBJECT INDEX: notation:expectations
\index{operators!collision expectation}
\index{operators!weighted range expectation}
\index{collision!probability}
% END SUBJECT INDEX

Our convention is
\[
 \boxed{N=\mathbb E_\lambda[KK^{\mathsf T}],\qquad
 M=\mathbb E_\lambda[K^{\mathsf T}JK].}
\]
Here the $K$ inside the expectation denotes a random kernel transformation,
and $J=\mathbf1\mathbf1^{\mathsf T}$.
Set $C_i=e_{ij}e_{ij}^{\mathsf T}$, independently of $j$. Its entries are
one precisely when their indices lie in the same block of $\mathcal P_i$.
Then
\[
 N=\sum_i\alpha_i C_i,\qquad N_{xy}=\Pr\{xK=yK\}.
\]
For $b(K)=\mathbf1^{\mathsf T}K$, the entries of this row are the fiber
sizes at the image vertices, and
\[
 M=\mathbb E[b(K)^{\mathsf T}b(K)].
\]
Thus $N$ describes partitions and $M$ describes weighted range data.
Both are symmetric positive semidefinite, but neither is generally a
projection. The pair-survival matrix is $J-N$, not another operator
named $N$. The rank of the average $N$ need not equal the kernel rank $r$.

\subsection{Column projections: hold the range fixed}
\label{notation:range-projections}
% BEGIN SUBJECT INDEX: notation:range-projections
\index{projection!range-fixed}
\index{kernel!range projections}
% END SUBJECT INDEX

Put $\rho_j=\operatorname{diag}(\mathbf1_{\mathcal R_j})$ and define
\[
 P_j=\sum_i\alpha_i e_{ij}.
\]
Idempotents in a fixed kernel column form a left-zero semigroup:
$e_{ij}e_{\ell j}=e_{ij}$. Therefore
\[
 P_j^2=\sum_{i,\ell}\alpha_i\alpha_\ell e_{ij}=P_j.
\]
Since $C_i\rho_j=e_{ij}$, we obtain the useful reconstruction formula
\[
 \boxed{P_j=N\rho_j.}
\]
On row vectors, $P_j$ projects onto vectors supported on $\mathcal R_j$:
its output is supported there and it fixes every coordinate row on that
range. In particular, $\operatorname{rank}P_j=r$.

\subsection{Row projections: hold the partition fixed}
\label{notation:partition-projections}
% BEGIN SUBJECT INDEX: notation:partition-projections
\index{projection!partition-fixed}
\index{kernel!partition projections}
\index{weighted block average}
% END SUBJECT INDEX

Define
\[
 Q_i=\sum_j\beta_j e_{ij}.
\]
Idempotents in a fixed kernel row form a right-zero semigroup:
$e_{ij}e_{i\ell}=e_{i\ell}$, so $Q_i^2=Q_i$.
For formulas involving the stationary row, assume the range-marginal
identity supplied by equipartition in the irreducible stochastic setting:
\[
 \sum_j\beta_j\rho_j=r\omega,\qquad
 \omega=\operatorname{diag}(\pi).
\]
This is a stated probabilistic hypothesis here, not a blanket assumption
on all finite semigroups. The PtC application must justify it in its own
setting when using these formulas. Averaging $C_i\rho_j=e_{ij}$ gives
\[
 \boxed{Q_i=rC_i\omega=r e_i e_i^{\mathsf T}\omega,}
\]
where $e_i$ is any idempotent with partition $\mathcal P_i$.
On column functions, $Q_i$ projects onto functions constant on each block:
\[
 (Q_i f)_x=r\sum_{y\in C}\pi_y f_y\quad(x\in C\in\mathcal P_i).
\]
Equipartition gives $\pi(C)=1/r$, making this a weighted block average.
It is not generally the Euclidean block-average projection. These
partition projections and the range projections above should also be
kept distinct from the higher-level limiting projections $\Omega_k$.

\subsection{The average of the distinguished idempotents}
\label{notation:average-idempotents}
% BEGIN SUBJECT INDEX: notation:average-idempotents
\index{idempotent!average of distinguished idempotents}
% END SUBJECT INDEX

Averaging the distinguished idempotents with the product marginal law gives
\[
 \boxed{\overline e=\sum_{i,j}\alpha_i\beta_j e_{ij}
 =\sum_i\alpha_iQ_i=\sum_j\beta_jP_j=rN\omega.}
\]
Although every $P_j$ and $Q_i$ is idempotent, this double average is not
asserted to be idempotent in general. These identities retain the actual
kernel $e$-projections; they are not replacements by unrelated orthogonal
projections.

